\chapter{Relation to established results and the remaining contribution}
\label{ch:rf-positioning}

The useful question is not whether partial observations can hide a mechanism.
That obstruction is established. The question pursued here is what a finite,
calibrated temporal record forces \emph{every compatible model} to contain or
expend. The answer depends on the experimental contract: a common dynamics,
repeatable preparation, detector, accessible times and an explicitly stated
classical or quantum comparison class. These conditions belong to the theorem,
not merely to its interpretation.

The work combines three tasks that are often separated: exclusion of models
with too few states, exclusion of stationary equilibrium explanations, and
constructive examples showing when neither conclusion is justified. The
specific mathematical and computational results should be evaluated on their
own statements. A shared presentation does not establish a new physical law,
a complete classification of hidden dynamics, or priority over the literatures
below.

\section{State number, positive realization and finite records}

Wolf and P\'erez-Garc\'ia already derived dimension information from observable
dynamics, including noisy Hankel-rank witnesses and a damped-qubit example
with arbitrarily large classical dimension requirements \cite{rrWolf2009}.
Their spectral argument uses stochastic eigenvalue restrictions originating
in work including Dmitriev--Dynkin and Karpelevich
\cite{rfDmitriev1946,rfKarpelevich1951}. An unbounded classical--qubit gap or a
frequency-to-decay state bound is therefore established prior work. Statistical
detection matrices also infer observable stochastic dimension from noisy data
\cite{rfSuccar2024}.

The extra issue is \emph{positive realization}: a scalar signal can have small
ordinary linear rank while requiring a much larger probabilistic realization.
The established positive-realization literature makes this distinction
explicit \cite{rfBenvenuti2022}. A nonnegative matrix factorization is a
necessary condition for a finite temporal model; it need not supply a common
stochastic semigroup realizing every time. Conversely, a low ordinary rank
cannot certify a small positive realization.

Two finite-record constructions address this gap here. The random-clock
construction supplies a quantitative distance from all generators below a
specified state number, allowing unknown preparation, unknown fixed detector
and arbitrarily large hidden rates. It does not infer an exact pole. The
antichain construction uses a finite temporal probability matrix and an
explicit positive-noise margin. The conservative record count of the first
construction is part of its result; existence of a finite test does not imply
experimental efficiency.

The antichain ingredient has especially close precedents. The exact
support-based inverse-central-binomial bound follows from established Boolean
rank and nonnegative-rank arguments \cite{rfDeCaen1981,rfGillis2020} and has
already been applied to quantum preparation--measurement tables
\cite{rfShahandeh2024}. Heinosaari and collaborators use the same sinusoidal
matrix, up to normalization, to demonstrate unbounded quantum advantage; they
also establish a stronger exact rank for the seven-point instance
\cite{rfHeinosaari2024}. The candidate addition here is the analytic tolerance
for nonzero diagonal entries, its single repeated-channel temporal realization,
and a specified finite-count rejection rule. It is not the first sinusoidal
matrix, quantum advantage or application of Sperner's theorem.

Temporal memory and environment-dimension witnesses provide another relevant
body of work \cite{rrVieira2022,rfVieira2024}. Their measurement sequences can
be invasive. A fresh preparation with one endpoint measurement imposes a
different contract, so numerical comparisons require matching the operations,
reset assumptions and counted memory. Quantum counting-process precision is
also connected to finite memory \cite{rfOlamaei2026}; an external randomized
endpoint clock is not itself a first-tick counting process.

\section{Classical approximation identifies the limit of the contract}

The Poisson smoothing operator used by the classical compiler is the
established Sz\'asz--Mirakjan construction \cite{rfSzasz1950}. Its approximation
properties are not new. The role of the compiler is to expose the resources
that can remain unobserved: enough hidden states, fast transitions and a
specially prepared initial state can approximate a finite collection of
prepare--wait--read curves over a bounded interval. An explicit finite Markov
implementation, error budget and adaptive-record comparison make that
obstruction usable in the present inference problem.

This yields a necessary distinction. A finite experiment can exclude every
model below a state bound while failing to exclude the union over all state
bounds. It can also reveal an expensive preparation without forcing persistent
stationary dissipation. A claim that the same finite trace disproves all
classical explanations would erase precisely the degrees of freedom that the
compiler constructs.

\section{Dissipation, reversible comparison and preparation}

Thermodynamic restrictions on coherence and temporal correlations are
established. They include affinity--coherence constraints, cross-correlation
asymmetry bounds, oscillation bounds and correlation-time inequalities
\cite{rfBarato2017,rfOhga2023,rfShiraishi2023,rfDechant2023,rfKolchinsky2024}.
Their resource definitions matter: cycle affinity, stationary entropy
production, activity and an eigenmode's weight are different quantities.
Recent work further emphasizes these distinctions
\cite{rfGu2026,rfShiraishi2026}. Finite-lag cross-correlation bounds are also
under active development \cite{rfRaux2026}.

The reversible comparison used here is not a new operation. Additive
reversibilization has an established information-geometric interpretation
\cite{rfChoi2024}, and the Jensen--Shannon/Jeffreys inequality underlying the
quarter-factor estimate is established \cite{rfCrooks2011}. The inspected,
unrefereed October 2026 manuscript \cite{rfHidden2026} already states the exact
trajectory bound
\[
 D(\mathbb P_Q^{[0,t]}\Vert\mathbb P_{(Q+Q^*)/2}^{[0,t]})
 \le t\sigma/4
\]
and uses reversible persistence restrictions to derive an
architecture-independent dissipation certificate. Its uninterrupted residence
observations differ from endpoint equality observations here. Its existence
precludes presenting the comparison inequality as an original discovery;
its unrefereed status means it does not independently validate this work.

The operational contribution pursued here uses the common positive spectral
measure of several endpoint autocorrelations. It certifies distance from that
moment set while propagating finite counts and detector calibration into an
entropy-production lower bound. A dual polynomial supplies a checkable lower
certificate, which a numerical best fit alone cannot provide. Binary
endpoint-pair records can be time-reversal symmetric even when their dependence
on lag violates the reversible spectral restrictions. The inference concerns
compatibility with a specified equilibrium model class, rather than an observed
pairwise current.

A further preparation--maintenance extension uses the standard adiabatic and
nonadiabatic decomposition of entropy production \cite{rfEsposito2010}. Initial
relative entropy and subsequent housekeeping dissipation are already known
physical resources. The candidate result is their combined lower constraint
from an accessible departure from equilibrium temporal laws. No claim is made
that the decomposition, initial nonequilibrium free energy, or data processing
is newly discovered.

Randomized observation times themselves also have close prior art. A recent
unrefereed manuscript uses two randomized clocks to identify fully observed
generators and studies the sampling cost of rejecting hidden equilibrium
explanations \cite{rfAznagulov2026}. This strengthens the need to state the
specific finite-record certificate, comparison class and error guarantee,
rather than claim first use of random clocks for hidden dynamics.

\section{What has and has not been established}

This part provides theorem statements, proofs and reproducible synthetic
examples. Independent mathematical checks and exact arithmetic support their
internal validity. The candidate additions are the particular robust temporal
certificates, their resource consequences and the construction of matching
obstructions under weaker experimental contracts. The literature audit
identifies their closest known ingredients and overlapping results. It does
not prove historical originality.

Applications to a physical apparatus require evidence for homogeneous dynamics,
preparation, readout, stationarity where invoked, and the relevant time-reversal
convention. In particular, even-state Markov entropy bounds do not apply without
modification to equilibrium systems with odd momentum variables or invasive
quantum measurements. No new empirical data or independent peer review is
reported. A practical next application should test an explicit laboratory
contract and compare against an existing method on the same records, rather
than infer engineering solutions from the breadth of the formalism.
